\SourcePage{263}
\SGAchapter{VI}{$\ell$ 进上同调}{Cohomologie $\ell$-adique}
\begin{center}J. P. Jouanolou\end{center}
本讲始终给定一个素数 $\ell$。与通常一样，以 $X_{\text{ét}}$ 表示\Term{prescheme}{预概形} $X$ 的\Term{etale-site}{平展景}。
\SGAsection{1}{可构造 $\ell$ 进层}
本节中，固定一个\Term{locally-noetherian}{局部诺特}\Term{prescheme}{预概形} $X$。
\SGAsection{1.1}{定义和一般性质}
\begin{SGAresult}{定义}{1.1.1}
所谓 $X$ 上的\Term{constructible-l-adic-sheaf}{可构造 $\ell$ 进层}，是指 $X_{\text{ét}}$ 上由\Term{constructible-sheaf}{可构造}\Term{abelian-sheaf}{阿贝尔层}组成的一个 $(\ell\mathbb Z)$ \Term{j-adic-system}{进投射系}（V 3.2.1）。
\end{SGAresult}
要使 $X_{\text{ét}}$ 上\Term{abelian-sheaf}{阿贝尔层}的一个 $(\ell\mathbb Z)$ \Term{j-adic-system}{进投射系} $F=(F_n)_{n\in\mathbb N}$ 成为\Term{constructible-l-adic-sheaf}{可构造 $\ell$ 进层}，只须 $F_0$ \Term{constructible}{可构造}。事实上，此断言是\Term{local-scheme}{局部}的，在\Term{noetherian}{诺特}情形由（V 5.1.2）得到。
以 $\mathrm{Abc}(X)$ 表示 $X$ 上\Term{constructible-sheaf}{可构造}\Term{abelian-sheaf}{阿贝尔层}的\Term{category}{范畴}。记
\[
\ell\text{-}\mathrm{adc}(X)
\]
为 $\operatorname{Hom}(\mathbb N^\circ,\mathrm{Abc}(X))$ 中由\Term{constructible-l-adic-sheaf}{可构造 $\ell$ 进层}生成的\Term{full-subcategory}{满子范畴}。
\SourcePage{264}
用第 V 讲的记号，有
\[
\ell\text{-}\mathrm{adc}(X)=(\ell\mathbb Z)\text{-}\mathrm{ad}(\mathrm{Abc}(X)).
\]
以 $\mathrm{Ab}(X)$ 表示 $X_{\text{ét}}$ 上\Term{abelian-sheaf}{阿贝尔层}的\Term{category}{范畴}。我们知道（SGA 4 IX 2.9），当 $X$ \Term{noetherian}{诺特}时，$X$ 上的\Term{constructible-sheaf}{可构造}\Term{abelian-sheaf}{阿贝尔层}恰是\Term{noetherian}{诺特}\Term{abelian-sheaf}{阿贝尔层}。此时，\Term{category}{范畴} $\ell\text{-}\mathrm{adc}(X)$ 显然可与第 V 讲中记作
\[
(\ell\mathbb Z)\text{-}\mathrm{ad}^{\mathrm n}(\mathrm{Ab}(X))
\]
的\Term{category}{范畴}等同。
\begin{SGAresult}{命题}{1.1.2}
设 $F=(F_n)_{n\in\mathbb N}$ 是\Term{category}{范畴} $\operatorname{Hom}(\mathbb N^\circ,\mathrm{Ab}(X))$ 的\Term{object}{对象}。$F$ 是\Term{constructible-l-adic-sheaf}{可构造 $\ell$ 进层}，当且仅当它对于 $X$ 的\Term{etale-topology}{平展拓扑}\Term{local-scheme}{局部}地具有此性质。
\end{SGAresult}
\begin{proof}
显然，$F$ 为 $\ell$ 进的性质可在此\Term{site}{景}的\Term{final-object}{终对象}的一个\Term{covering}{覆盖}上判定。另一方面，由（SGA 4 IX 2.8），若各 $F_n$ \Term{local-scheme}{局部}\Term{constructible}{可构造}，则它们就是\Term{constructible}{可构造}的。
\end{proof}
\begin{SGAresult}{命题}{1.1.3}
$X$ 的\Term{constructible-l-adic-sheaf}{可构造 $\ell$ 进层}\Term{category}{范畴} $\ell\text{-}\mathrm{adc}(X)$ 是\Term{abelian-category}{阿贝尔范畴}。当 $X$ \Term{noetherian}{诺特}时，它也是\Term{noetherian}{诺特}的。
\end{SGAresult}
\begin{proof}
当 $X$ \Term{noetherian}{诺特}时，这两个断言都由（V 5.2.3）得到。一般情形下，需要证明 $X$ 上\Term{constructible-l-adic-sheaf}{可构造 $\ell$ 进层}的任一\Term{morphism}{态射}
\[
u:F\longrightarrow G
\]
都有\Term{kernel}{核}和\Term{cokernel}{余核}。
\SourcePage{265}
\Term{cokernel}{余核}的存在性由（V 3.1.3）得到。现证 $u$ 的\Term{kernel}{核}存在。对 $X$ 的每个\Term{noetherian}{诺特}\Term{open-subset}{开集} $V$，我们知道 $u$ 在 $V$ 上的\Term{restriction}{限制}
\[
u_V:F_V\longrightarrow G_V
\]
在\Term{category}{范畴} $\ell\text{-}\mathrm{adc}(V)$ 中有\Term{kernel}{核}，而且有如下简单描述。若 $K$ 是 $u_V$ 在\Term{category}{范畴} $\operatorname{Hom}(\mathbb N^\circ,\mathrm{Ab}(V))$ 中的\Term{kernel}{核}，则 $K$ 满足 \Term{mlar-condition}{Mittag--Leffler--Artin--Rees 条件}；以 $K'$ 表示它的\Term{universal-images}{泛像投射系}，则存在一个整数 $r$，使\Term{projective-system}{投射系}
\[
\mathrm m_r(K')=\bigl(K'_{n+r}/\ell^{n+1}K'_{n+r}\bigr)_{n\in\mathbb N}
\]
（取显然的\Term{transition-morphism}{过渡态射}）为\Term{constructible-l-adic-sheaf}{可构造 $\ell$ 进层}。下面显然的复合在\Term{category}{范畴} $\ell\text{-}\mathrm{adc}(V)$ 中是 $u_V$ 的一个\Term{kernel}{核}：
\[
\mathrm m_r(K')\longrightarrow K'\longrightarrow K.
\]
上述构造特别表明，若 $W$ 是包含于 $V$ 的一个\Term{open-subset}{开集}，则 $u_V$ 的一个\Term{kernel}{核}在 $W$ 上的\Term{restriction}{限制}，是 $u_W$ 的一个\Term{kernel}{核}。
于是，设 $(V_i)_{i\in I}$ 是 $X$ 的一个由\Term{noetherian}{诺特}\Term{open-subset}{开集}组成的\Term{covering}{覆盖}。若 $i,j\in I$，则 $u_{V_i}$ 和 $u_{V_j}$ 的\Term{kernel}{核}在 $V_i\cap V_j$ 上的\Term{restriction}{限制}\Term{canonical-isomorphism}{典范同构}；把它们\Term{gluing}{粘合}起来，就得到 $u$ 的一个\Term{kernel}{核}。
\end{proof}
\textbf{1.1.4。} 设 $F=(F_n)_{n\in\mathbb N}$ 是\Term{constructible-l-adic-sheaf}{可构造 $\ell$ 进层}。环 $\mathbb Z_\ell$ 显然\Term{action}{作用}于每个\Term{sheaf}{层} $F_n$，因而也\Term{action}{作用}于 $F$。此外，若 $F,G$ 是两个\Term{constructible-l-adic-sheaf}{可构造 $\ell$ 进层}，且
\[
u:F\longrightarrow G
\]
\SourcePage{266}
是一个\Term{morphism}{态射}，则 $u$ 与 $\mathbb Z_\ell$ 分别在 $F$ 和 $G$ 上的\Term{action}{作用}相容。这表明 $\ell\text{-}\mathrm{adc}(X)$ \Term{canonical}{典范}地带有一个\Term{abelian-a-category}{阿贝尔 $\mathbb Z_\ell$ 范畴}的结构。
因此可以采用如下用语。

\textbf{用语。} 下文把\Term{constructible-l-adic-sheaf}{可构造 $\ell$ 进层}称为\Term{constructible-zl-sheaf}{可构造 $\mathbb Z_\ell$ 层}。
须注意，不要把这一概念与 $(\mathbb Z_\ell)_X$ 模的概念混淆；下文不使用后一个概念。
\Term{category}{范畴} $\ell\text{-}\mathrm{adc}(X)$ 也记作
\[
\mathbb Z_\ell\text{-}\mathrm{fc}(X).
\]
\SGAsection{1.2}{可构造扭常值 $\mathbb Z_\ell$ 层}
\begin{SGAresult}{定义}{1.2.1}
一个\Term{constructible-zl-sheaf}{可构造 $\mathbb Z_\ell$ 层}
\[
F=(F_n)_{n\in\mathbb N}
\]
称为\Term{twisted-constant}{扭常值}（constant tordu）的，当且仅当各\Term{sheaf}{层} $F_n$ 在 $X_{\text{ét}}$ 上\Term{locally-constant}{局部常值}。
由\Term{constructible-zl-sheaf}{可构造}\Term{twisted-constant}{扭常值 $\mathbb Z_\ell$ 层}在 $\mathbb Z_\ell\text{-}\mathrm{fc}(X)$ 中生成的\Term{full-subcategory}{满子范畴}，称为\Term{constructible-zl-sheaf}{可构造}\Term{twisted-constant}{扭常值 $\mathbb Z_\ell$ 层}\Term{category}{范畴}。
\end{SGAresult}
\textbf{例 1.2.2。$\mathbb Z_\ell$ 层 $\mathbb Z_\ell(1)$。}
假设整数 $\ell$ 与 $X$ 的各\Term{residue-characteristic}{剩余特征}互素。对每个整数 $n$，取 $\ell$ 次幂给出\Term{locally-constant}{局部常值}\Term{sheaf}{层}的一个\Term{morphism}{态射}：
\SourcePage{267}
\[
\mu_{\ell^{n+1}}\longrightarrow\mu_{\ell^n}.
\]
这显然定义了一个\Term{projective-system}{投射系}：
\[
\mathbb Z_\ell(1)=(\mu_{\ell^{n+1}})_{n\in\mathbb N}.
\]
对每个整数 $n$，归约到 $\mu_{\ell^n}$ 和 $\mu_{\ell^{n+1}}$ 都是\Term{constant-sheaf}{常值}的情形，可知\Term{projective-system}{投射系} $\mathbb Z_\ell(1)$ 是 $\ell$ 进的。
当 $X=\operatorname{Spec}(K)$，其中 $K$ 是其\Term{prime-field}{素域}的\Term{finite-type}{有限型}\Term{group-extension}{扩张}时，$\mathbb Z_\ell$ 层 $\mathbb Z_\ell(1)$ 给出一个\Term{constructible-l-adic-sheaf}{可构造}\Term{twisted-constant}{扭常值 $\ell$ 进层}（1.2.1）的例子，使得不存在 $K$ 的\Term{finite-etale}{有限平展}\Term{group-extension}{扩张} $K'$，令各 $F_n$ 在 $X'=\operatorname{Spec}(K')$ 上的\Term{inverse-image}{逆像}都为\Term{constant-sheaf}{常值}。否则，对每个整数 $n$，$K'$ 都将包含一个本原 $\ell^n$ 次\Term{root-of-unity}{单位根}，这对\Term{prime-field}{素域}的\Term{finite-type}{有限型}\Term{group-extension}{扩张}是不可能的。

\textbf{回顾 1.2.3（SGA 1 V 7）。} 我们知道，\Term{prescheme}{预概形} $T$ 上\Term{locally-constant}{局部常值}的\Term{constructible-sheaf}{可构造}\Term{torsion}{挠}\Term{abelian-sheaf}{阿贝尔层}\Term{category}{范畴}，\Term{equivalence}{等价}于带有阿贝尔 $X$ \Term{group-structure}{群结构}的 $T$ 的\Term{etale-covering}{平展覆盖}\Term{category}{范畴}。现设 $X$ 是一个\Term{connected}{连通}的\Term{locally-noetherian}{局部诺特}\Term{prescheme}{预概形}。以 $\mathrm R_X$ 表示 $X$ 的\Term{etale-covering}{平展覆盖}\Term{category}{范畴}，并以 $\mathrm{Ensf}$ 表示\Term{finite}{有限}集\Term{category}{范畴}。若 $a$ 是 $X$ 的一个\Term{geometric-point}{几何点}，则令
\[
F_a:\mathrm R_X\longrightarrow\mathrm{Ensf}
\]
\SourcePage{268}
为把 $X$ 的每个\Term{etale-covering}{平展覆盖}对应到其在点 $a$ 的\Term{geometric-fiber}{几何纤维}的\Term{functor}{函子}。按定义，$X$ 在点 $a$ 的\Term{fundamental-group}{基本群}
\[
\pi_1(X)=\pi_1(X,a)
\]
是 $F_a$ 的\Term{automorphism-group}{自同构群}，配以适当的\Term{topology}{拓扑}。因此 $F_a$ 可以解释为一个\Term{functor}{函子}
\[
\mathrm R_X\longrightarrow\mathrm{Ensf}(\pi_1),
\]
从\Term{category}{范畴} $\mathrm R_X$ 映到带有 $\pi_1=\pi_1(X,a)$ 的\Term{continuous-action}{连续作用}的\Term{finite}{有限}集\Term{category}{范畴}。
回顾如下命题。
\begin{SGAresult}{命题}{1.2.3.1}
\Term{functor}{函子}
\[
F_a:\mathrm R_X\longrightarrow\mathrm{Ensf}(\pi_1)
\]
是一个\Term{equivalence-of-categories}{范畴等价}。它与\Term{inductive-limit}{归纳极限}及\Term{finite}{有限}\Term{projective-limit}{投射极限}交换。
\end{SGAresult}
现在分别以
\[
\mathrm{Lc}(X),\qquad\mathrm{Abf}(\pi_1)
\]
表示\Term{locally-constant}{局部常值}的\Term{constructible-sheaf}{可构造}\Term{torsion}{挠}\Term{abelian-sheaf}{阿贝尔层}\Term{category}{范畴}，以及带有\Term{fundamental-group}{基本群}的\Term{continuous-action}{连续作用}的\Term{finite}{有限}\Term{abelian-group}{阿贝尔群}\Term{category}{范畴}。\Term{functor}{函子} $F_a$ 诱导一个仍记作 $F_a$ 的\Term{exact-functor}{正合函子}
\[
F_a:\mathrm{Lc}(X)\longrightarrow\mathrm{Abf}(\pi_1),
\]
它是一个\Term{equivalence-of-categories}{范畴等价}。
\SourcePage{269}
\textbf{1.2.4。\Term{connected}{连通}\Term{prescheme}{预概形}上的\Term{twisted-constant}{扭常值} $\mathbb Z_\ell$ 层。}
由（1.2.3）的回顾，立即得到下面的\Term{lemma}{引理}。
\begin{SGAresult}{\Term{lemma}{引理}}{1.2.4.1}
把 $F_a$ 延拓到由 $\mathbb N$ 索引的\Term{projective-system}{投射系}\Term{category}{范畴}，所得\Term{functor}{函子} $\widehat F_a$ 诱导\Term{constructible-l-adic-sheaf}{可构造}\Term{twisted-constant}{扭常值 $\ell$ 进层}\Term{category}{范畴}与\Term{category}{范畴}
\[
(\ell\mathbb Z)\text{-}\mathrm{ad}(\mathrm{Abf}(\pi_1))
\]
之间的一个\Term{equivalence-of-categories}{等价}。
\end{SGAresult}
现设 $F=(F_n)_{n\in\mathbb N}$ 是\Term{category}{范畴} $(\ell\mathbb Z)\text{-}\mathrm{ad}(\mathrm{Abf}(\pi_1))$ 的\Term{object}{对象}。它的\Term{projective-limit}{投射极限} $\varprojlim_{\mathbb N}(F_n)$ 自然带有\Term{finite-type}{有限型} $\mathbb Z_\ell$ 模结构，$\pi_1$ 在其上对于 $\ell$ \Term{adic-topology}{进拓扑}\Term{continuous-action}{连续作用}。回顾，若 $G$ 和 $M$ 分别表示一个\Term{profinite-group}{投有限群}和一个\Term{finite-type}{有限型} $\mathbb Z_\ell$ 模，则所谓 $G$ 在 $M$ 上的\Term{continuous-action}{连续作用}，是指一个\Term{continuous-map}{连续映射}
\[
G\times M\longrightarrow M,\qquad(g,m)\longmapsto g.m,
\]
满足
\[
\begin{aligned}
\text{(i)}\quad&g.(m+n)=g.m+g.n&&(g\in G,\ m,n\in M),\\
\text{(ii)}\quad&(gh).m=g.(h.m)&&(g,h\in G,\ m\in M).
\end{aligned}
\]
于是，显然对每个 $\ell$ 进整数 $a$，都有
\[
g.(am)=a(g.m)\qquad(g\in G,\ m\in M).
\]
另一方面，以 $\mathrm{GL}(M)$ 表示从 $M$ 到自身的\Term{invertible}{可逆} $\mathbb Z_\ell$ \Term{linear}{线性}映射组成的群，它可与各群 $\mathrm{GL}(M/\ell^{n+1}M)$ 的\Term{projective-limit}{投射极限}等同。
\SourcePage{270}
这些群都是\Term{finite-group}{有限群}，因而群 $\mathrm{GL}(M)$ 自然带有\Term{profinite-group}{投有限群}结构，特别带有紧\Term{topological-group}{拓扑群}结构。容易验证，给定\Term{profinite-group}{投有限群} $G$ 在 $M$ 上的一个\Term{continuous-action}{连续作用}，\Term{equivalence}{等价}于给定一个\Term{topological-group}{拓扑群}\Term{morphism}{态射}
\[
G\longrightarrow\mathrm{GL}(M),
\]
其中后一个群取上述\Term{topology}{拓扑}。
由此，显然可定义一个\Term{functor}{函子}
\[
\varprojlim_{\mathbb N}:(\ell\mathbb Z)\text{-}\mathrm{ad}(\mathrm{Abf}(\pi_1))
\longrightarrow\mathbb Z_\ell\text{-}\mathrm{modn}(\pi_1),
\]
其中 $\mathbb Z_\ell\text{-}\mathrm{modn}(\pi_1)$ 表示带有 $\pi_1$ 的\Term{continuous-action}{连续作用}的\Term{finite-type}{有限型} $\mathbb Z_\ell$ \Term{module-category}{模范畴}。这使我们可以陈述下述\Term{lemma}{引理}。
\begin{SGAresult}{\Term{lemma}{引理}}{1.2.4.2}
上述\Term{functor}{函子}
\[
\varprojlim_{\mathbb N}:(\ell\mathbb Z)\text{-}\mathrm{ad}(\mathrm{Abf}(\pi_1))
\longrightarrow\mathbb Z_\ell\text{-}\mathrm{modn}(\pi_1)
\]
是一个\Term{equivalence-of-categories}{范畴等价}。
\end{SGAresult}
\begin{proof}
显然，它在\Term{object}{对象}上\Term{essentially-surjective}{本质满射}：一个带有 $\pi_1$ 的\Term{continuous-action}{连续作用}的\Term{finite-type}{有限型} $\mathbb Z_\ell$ 模 $E$，是\Term{projective-system}{投射系}
\[
(E/\ell^{n+1}E)_{n\in\mathbb N}
\]
的像，各项取显然的 $\pi_1$ \Term{action}{作用}。现设 $F=(F_n)_{n\in\mathbb N}$ 和 $G=(G_n)_{n\in\mathbb N}$ 是\Term{category}{范畴} $(\ell\mathbb Z)\text{-}\mathrm{ad}(\mathrm{Abf}(\pi_1))$ 的两个\Term{object}{对象}；为简便起见，以 $\operatorname{Hom}_{\pi_1}(F,G)$ 表示该\Term{category}{范畴}中从 $F$ 到 $G$ 的\Term{morphism}{态射}集。
\SourcePage{271}
忘掉 $\pi_1$ 的\Term{action}{作用}，可以把 $F$ 和 $G$ 看成\Term{finite}{有限}\Term{abelian-group}{阿贝尔群}的 $\ell$ 进\Term{projective-system}{投射系}\Term{category}{范畴}中的\Term{object}{对象}，并以 $\operatorname{Hom}_0(F,G)$ 表示该\Term{category}{范畴}中从 $F$ 到 $G$ 的\Term{morphism}{态射}集。群 $\pi_1$ 如下\Term{action}{作用}于 $\operatorname{Hom}_0(F,G)$：$\pi_1$ 的一个元素 $s$ 显然分别定义 $F$ 和 $G$ 在\Term{finite}{有限}\Term{abelian-group}{阿贝尔群}的 $\ell$ 进\Term{projective-system}{投射系}\Term{category}{范畴}中的\Term{automorphism}{自同构}，均记作 $s$；若 $u\in\operatorname{Hom}_0(F,G)$，则置
\[
s.u=s\circ u\circ s^{-1}.
\]
于是，$\operatorname{Hom}_{\pi_1}(F,G)$ 显然\Term{canonical}{典范}地等同于 $\operatorname{Hom}_0(F,G)$ 中 $\pi_1$ 的\Term{invariant}{不变}元\Term{set}{集合}；同样，对每个 $n$，$\operatorname{Hom}_{\pi_1}(F_n,G_n)$ 等同于 $\operatorname{Hom}_{\mathbb Z}(F_n,G_n)^{\pi_1}$，其中 $\pi_1$ 的\Term{action}{作用}以类似方式定义。例如由 Serre：《Groupes proalgébriques》（《投\Term{algebraic-group}{代数群}》），命题 12，可知若以 $M,N$ 分别表示 $F,G$ 的\Term{projective-limit}{投射极限}，则
\[
\operatorname{Hom}_{\mathbb Z_\ell}(M,N)
=\varprojlim_{\mathbb N}\operatorname{Hom}_{\mathbb Z}(F_n,G_n).
\]
在两边取 $\pi_1$ \Term{invariant}{不变}元，并利用此运算对应于一个\Term{projective-limit}{投射极限}，故与\Term{projective-limit}{投射极限}交换，即得到所需断言。
\end{proof}
\Term{lemma}{引理} 1.2.4.1 和 1.2.4.2 给出下述较不精细的命题。
\begin{SGAresult}{命题}{1.2.5}
一个\Term{connected}{连通}的\Term{locally-noetherian}{局部诺特}\Term{prescheme}{预概形} $X$ 上\Term{constructible-zl-sheaf}{可构造}\Term{twisted-constant}{扭常值 $\mathbb Z_\ell$ 层}的\Term{category}{范畴}，\Term{equivalence}{等价}于 $\pi_1(X)$ 在其上对于 $\ell$ \Term{adic-topology}{进拓扑}\Term{continuous-action}{连续作用}的\Term{finite-type}{有限型} $\mathbb Z_\ell$ \Term{module-category}{模范畴}。
\end{SGAresult}
\SourcePage{272}
\textbf{注记。} 设 $F$ 是\Term{constructible-zl-sheaf}{可构造}\Term{twisted-constant}{扭常值 $\mathbb Z_\ell$ 层}，$M$ 是所对应的、带有 $\pi_1$ 的\Term{continuous-action}{连续作用}的\Term{finite-type}{有限型} $\mathbb Z_\ell$ 模。要使 $F$ \Term{locally-constant}{局部常值}，显然只须 $\pi_1$ 在 $M$ 上的\Term{action}{作用}通过一个\Term{finite}{有限}\Term{quotient-group}{商群}\Term{resolution}{分解}，而这一条件并不总成立（1.2.2）。当 $X$ \Term{normal}{正规}时，这个条件也是必要的。把 $X$ \Term{resolution}{分解}为\Term{connected-component}{连通分支}，即可归约到 $X$ \Term{integral}{整}的情形。以 $K$ 表示它的\Term{rational-function-field}{有理函数域}；我们断言，若存在 $K$ 的\Term{finite-separable-extension}{有限可分扩张} $K'$，使 $F$ 在 $\operatorname{Spec}(K')$ 上的\Term{inverse-image}{逆像} $F_{K'}$ 为\Term{constant-sheaf}{常值}，则群 $\pi_1(X)$ 在 $M$ 上的\Term{action}{作用}通过一个\Term{finite-group}{有限群}\Term{resolution}{分解}。事实上（SGA 1 V 8.2），$\pi_1(K)\to\pi_1(X)$ 是\Term{surjectivity}{满射}，而假设意味着复合
\[
\pi_1(K)\longrightarrow\pi_1(X)\longrightarrow\operatorname{Aut}(M)
\]
的\Term{kernel}{核}具有\Term{finite-index}{有限指数}。因此，$\pi_1(X)\to\operatorname{Aut}(M)$ 的\Term{kernel}{核}也有\Term{finite-index}{有限指数}。
\begin{SGAresult}{命题}{1.2.6}
设 $X$ 是一个\Term{locally-noetherian}{局部诺特}\Term{prescheme}{预概形}，$F$ 是 $X_{\text{ét}}$ 上由 $\mathbb N$ 索引的\Term{abelian-sheaf}{阿贝尔层}\Term{projective-system}{投射系}。下列断言\Term{equivalence}{等价}：

(i) $F$ 是\Term{constructible-zl-sheaf}{可构造 $\mathbb Z_\ell$ 层}；

(ii) $X$ 的每个\Term{noetherian}{诺特}\Term{open-subset}{开集} $U$ 都是\Term{finite}{有限}个\Term{locally-closed}{局部闭}\Term{subset}{子集} $(Z_i)_{1\leq i\leq q}$ 的并，且 $F$ 在每个 $Z_i$ 上的\Term{inverse-image}{逆像}都是\Term{constructible-zl-sheaf}{可构造}\Term{twisted-constant}{扭常值 $\mathbb Z_\ell$ 层}；

(iii) $X$ 有一个开\Term{covering}{覆盖}，其中每个\Term{open-subset}{开集}都是\Term{finite}{有限}个\Term{locally-closed}{局部闭}\Term{subset}{子集}的并，且 $F$ 在这些\Term{subset}{子集}上的\Term{inverse-image}{逆像}都是\Term{constructible-zl-sheaf}{可构造}\Term{twisted-constant}{扭常值 $\mathbb Z_\ell$ 层}。
\end{SGAresult}
\SourcePage{273}
\begin{proof}
显然，(ii) 蕴含 (iii)。

(iii)$\Rightarrow$(i)：若 $F$ 满足 (iii)，则它的\Term{geometric-fiber}{几何纤维}都是 $\ell$ 进的，故由（SGA 4 VIII 3.6），$F$ 是 $\ell$ 进的。另一方面，它的各\Term{set-component}{分量}都是\Term{constructible}{可构造}的（SGA 4 IX 2.8）。

为证 (i) 蕴含 (ii)，立即可归约到 $X$ \Term{noetherian}{诺特}的情形。若 $F$ 是\Term{constructible-l-adic-sheaf}{可构造 $\ell$ 进层}，则它的\Term{strict-graded}{严格分次对象} $\operatorname{grs}(F)$ 是一个\Term{noetherian}{诺特} $(\operatorname{gr}\ell\mathbb Z,\mathrm{Ab}(X))$ 模（V 5.1.6）。因此（SGA 4 IX 2.9），$X$ 是\Term{finite}{有限}个\Term{locally-closed}{局部闭}\Term{subset}{子集} $(Z_i)_{1\leq i\leq q}$ 的并，且 $\operatorname{grs}(F)$ 在每个 $Z_i$ 上的\Term{inverse-image}{逆像}都是\Term{locally-constant}{局部常值}\Term{sheaf}{层}。由于具有\Term{finite}{有限}\Term{fiber}{纤维}的\Term{sheaf}{层}的\Term{locally-constant}{局部常值}性质在\Term{group-extension}{扩张}下\Term{stable}{稳定}，可\Term{pushout}{推出}，对每个整数 $n$，$F$ 的\Term{set-component}{分量} $F_n$ 在每个 $Z_i$ 上的\Term{inverse-image}{逆像}都是\Term{locally-constant}{局部常值}\Term{sheaf}{层}。
\end{proof}
\SGAsection{1.3}{$\mathbb Z_\ell$ 层上的 $\otimes$ 与 $\mathcal{H}om$ 运算}
\textbf{1.3.1。\Term{tensor-product}{张量积}。} 设 $F,G$ 是\Term{category}{范畴} $\operatorname{Hom}(\mathbb N^\circ,\mathrm{Ab}(X))$ 的两个\Term{object}{对象}。若 $F=(F_n)_{n\in\mathbb N}$、$G=(G_n)_{n\in\mathbb N}$，则在\Term{object}{对象}上用下式定义 $F$ 和 $G$ 的\Term{tensor-product}{张量积} $F\otimes G$，在\Term{arrow}{箭头}上的定义是显然的：
\[
(F\otimes G)_n=F_n\otimes G_n.
\]
更确切地说，这显然定义了一个\Term{additive}{加性}的\Term{right-exact-functor}{右正合}\Term{bifunctor}{双函子}，称为\Term{tensor-product}{张量积}：
\[
\otimes:\operatorname{Hom}(\mathbb N^\circ,\mathrm{Ab}(X))\times
\operatorname{Hom}(\mathbb N^\circ,\mathrm{Ab}(X))
\longrightarrow\operatorname{Hom}(\mathbb N^\circ,\mathrm{Ab}(X)).
\]
\SourcePage{274}
\begin{SGAresult}{命题}{1.3.2}
若 $F,G$ 是两个\Term{constructible-zl-sheaf}{可构造 $\mathbb Z_\ell$ 层}，则\Term{tensor-product}{张量积} $F\otimes G$ 是\Term{constructible-zl-sheaf}{可构造 $\mathbb Z_\ell$ 层}。
\end{SGAresult}
\begin{proof}
$\ell$ 进性是显然的。另一方面，由\Term{constructible-sheaf}{可构造层}的定义（SGA 4 IX 2.3），立即可知两个\Term{constructible-sheaf}{可构造层}的\Term{tensor-product}{张量积}是\Term{constructible}{可构造}的。
\end{proof}
因此，取\Term{restriction}{限制}可定义一个\Term{additive}{加性}的\Term{right-exact-functor}{右正合}\Term{bifunctor}{双函子}，仍称为\Term{tensor-product}{张量积}：
\[
\otimes:\mathbb Z_\ell\text{-}\mathrm{fc}(X)\times\mathbb Z_\ell\text{-}\mathrm{fc}(X)
\longrightarrow\mathbb Z_\ell\text{-}\mathrm{fc}(X).
\]
\textbf{1.3.3。\Term{hom-sheaf}{同态层}。}
设 $F=(F_n)_{n\in\mathbb N}$、$G=(G_n)_{n\in\mathbb N}$ 是两个\Term{constructible-zl-sheaf}{可构造 $\mathbb Z_\ell$ 层}。假设 $F$ \Term{locally-free}{局部自由}，意思是，对每个整数 $n$，\Term{sheaf}{层} $F_n$ 在\Term{constant-sheaf}{常值层} $(\mathbb Z/\ell^{n+1}\mathbb Z)_X$ 上\Term{locally-free}{局部自由}。对每个整数 $n$，结构\Term{exact-sequence}{正合列}
\[
G_{n+1}\xrightarrow{\ell^{n+1}}G_{n+1}\longrightarrow G_n\longrightarrow0
\]
给出下述\Term{exact-sequence}{正合列}，因为 $F_{n+1}$ 在 $(\mathbb Z/\ell^{n+1}\mathbb Z)_X$ 上\Term{locally-free}{局部自由}：
\[
\mathcal{H}om(F_{n+1},G_{n+1})\xrightarrow{\ell^{n+1}}
\mathcal{H}om(F_{n+1},G_{n+1})\longrightarrow
\mathcal{H}om(F_{n+1},G_n)\longrightarrow0.\tag{1}
\]
由于 $\ell^{n+1}G_n=0$，结构\Term{exact-sequence}{正合列}
\[
F_{n+1}\xrightarrow{\ell^{n+1}}F_{n+1}\longrightarrow F_n\longrightarrow0
\]
显然定义了一个\Term{isomorphism}{同构}：
\SourcePage{275}
\[
\mathcal{H}om(F_{n+1},G_n)\xrightarrow{\sim}\mathcal{H}om(F_n,G_n).\tag{2}
\]
结合 (1)、(2)，对每个整数 $n$ 得到一个\Term{morphism}{态射}
\[
\mathcal{H}om(F_{n+1},G_{n+1})\longrightarrow\mathcal{H}om(F_n,G_n).
\]
于是得到一个\Term{projective-system}{投射系}
\[
\mathcal{H}om(F,G)=\bigl(\mathcal{H}om(F_n,G_n)\bigr)_{n\in\mathbb N},
\]
它显然是一个 $\mathbb Z_\ell$ \Term{sheaf}{层}。

\textbf{1.3.4。\Term{invertible}{可逆} $\ell$ 进层的一般性质；在层 $\mathbb Z_\ell(r)$（$r\in\mathbb Z$）上的应用。}
设 $L$ 是一个\Term{constructible-zl-sheaf}{可构造 $\mathbb Z_\ell$ 层}。若对每个整数 $n$，它的第 $n$ 个\Term{set-component}{分量} $L_n$ 都是一个\Term{invertible-module}{可逆} $(\mathbb Z/\ell^{n+1}\mathbb Z)_X$ 模，即\Term{locally-free}{局部自由}秩为一的模，则称 $L$ \Term{invertible-l-adic-sheaf}{可逆}。这\Term{equivalence}{等价}于说，$L$ 在上一段的意义下\Term{locally-free}{局部自由}，并且各 $L_n$ 的\Term{geometric-fiber}{几何纤维}都是非零循环\Term{abelian-group}{阿贝尔群}。
若 $L$ 是一个\Term{invertible-l-adic-sheaf}{可逆 $\mathbb Z_\ell$ 层}，约定置
\[
L^{-1}=\mathcal{H}om(L,\mathbb Z_\ell).
\]
显然，存在\Term{isomorphism}{同构}
\[
L\otimes L^{-1}\xrightarrow{\sim}\mathbb Z_\ell,
\]
它由逐\Term{set-component}{分量}的相应\Term{isomorphism}{同构}得到（EGA $0_{\mathrm I}$ 5.4.3.1）。接着，依照（EGA $0_{\mathrm I}$ 5.4.4），约定置 $L^{\otimes0}=\mathbb Z_\ell$，并且当 $n\geq1$ 时置 $L^{\otimes(-n)}=(L^{-1})^{\otimes n}$。采用这些记号，对任意整数 $m,n$，便有一个\Term{canonical}{典范}的、\Term{functorial}{函子性}的\Term{isomorphism}{同构}
\[
L^{\otimes m}\otimes L^{\otimes n}\xrightarrow{\sim}L^{\otimes(m+n)}.
\]
\SourcePage{276}
此外，这一\Term{isomorphism}{同构}满足通常的\Term{associativity}{结合性}质。特别地，\Term{graded-object}{分次对象}
\[
(L^{\otimes i})_{i\in\mathbb N}\qquad\text{和}\qquad(L^{\otimes i})_{i\in\mathbb Z}
\]
是分别由 $\mathbb N$ 和 $\mathbb Z$ \Term{graded}{分次}的 $\mathbb Z_\ell$ \Term{algebra}{代数}。

\textbf{例。} （1.2.2）中定义的 $\mathbb Z_\ell$ \Term{sheaf}{层} $\mathbb Z_\ell(1)$ 是\Term{invertible-l-adic-sheaf}{可逆}的；下文约定置
\[
\mathbb Z_\ell(r)=(\mathbb Z_\ell(1))^{\otimes r}\qquad(r\in\mathbb Z).
\]
\SGAsection{1.4}{由可构造 $\mathbb Z_\ell$ 层范畴构造的各种范畴}
\textbf{1.4.1。\Term{constructible}{可构造} $A$ 层（$A$ 为\Term{finite}{有限} $\mathbb Z_\ell$ \Term{algebra}{代数}）。}
设 $A$ 是一个\Term{finite}{有限} $\mathbb Z_\ell$ \Term{algebra}{代数}。如下定义一个记作
\[
A\text{-}\mathrm{fc}(X)
\]
的\Term{category}{范畴}，其\Term{object}{对象}称为\Term{constructible-a-sheaf}{可构造 $A$ 层}。

(i) $A\text{-}\mathrm{fc}(X)$ 的一个\Term{object}{对象}是一个偶 $(F,u)$，其中 $F$ 为\Term{constructible-zl-sheaf}{可构造 $\mathbb Z_\ell$ 层}，$u$ 是一个 $\mathbb Z_\ell$ \Term{algebra}{代数}\Term{morphism}{态射}
\[
u:A\longrightarrow\operatorname{End}(F).
\]
(ii) 若 $(F,u)$、$(F',u')$ 是两个\Term{constructible-a-sheaf}{可构造 $A$ 层}，则一个\Term{arrow}{箭头} $(F,u)\to(F',u')$ 是一个 $\mathbb Z_\ell$ \Term{sheaf}{层}的\Term{morphism}{态射}
\[
f:F\longrightarrow F',
\]
使得对 $A$ 的每个元素 $a$，下图交换：
\SourcePage{277}
\[
\begin{tikzcd}
F\arrow[r,"f"]\arrow[d,"u(a)"']&F'\arrow[d,"u'(a)"]\\
F\arrow[r,"f"]&F'.
\end{tikzcd}
\]
一个\Term{constructible-a-sheaf}{可构造 $A$ 层}，若其底层的\Term{constructible-zl-sheaf}{可构造 $\mathbb Z_\ell$ 层}是\Term{twisted-constant}{扭常值}的，则称为\Term{constructible-a-sheaf}{可构造}\Term{twisted-constant}{扭常值 $A$ 层}。由\Term{constructible-a-sheaf}{可构造}\Term{twisted-constant}{扭常值 $A$ 层}在\Term{constructible-a-sheaf}{可构造 $A$ 层}\Term{category}{范畴}中生成的\Term{full-subcategory}{满子范畴}，称为\Term{constructible-a-sheaf}{可构造}\Term{twisted-constant}{扭常值 $A$ 层}\Term{category}{范畴}。
由（1.2.4）立即可知，若 $X$ \Term{connected}{连通}，则\Term{constructible-a-sheaf}{可构造}\Term{twisted-constant}{扭常值 $A$ 层}\Term{category-equivalence}{范畴等价}于以下\Term{category}{范畴}：其\Term{object}{对象}是\Term{finite-type}{有限型} $A$ 模，在底层的\Term{finite-type}{有限型} $\mathbb Z_\ell$ 模上带有 $\pi_1(X)$ 的\Term{continuous}{连续}且 $A$ \Term{linear}{线性}的\Term{action}{作用}（1.2.4）。

\textbf{1.4.2。\Term{constructible}{可构造} $\mathbb Q_\ell$ 层。} 显然，由\Term{constructible}{可构造}\Term{torsion}{挠} $\mathbb Z_\ell$ \Term{sheaf}{层}（即被 $\ell$ 的某个幂所\Term{annihilate}{零化}的\Term{sheaf}{层}）在 $\mathbb Z_\ell\text{-}\mathrm{fc}(X)$ 中生成的\Term{full-subcategory}{满子范畴}是\Term{thick-subcategory}{厚}的。
\begin{SGAresult}{定义}{1.4.3}
$X$ 上的\Term{constructible-ql-sheaf}{可构造 $\mathbb Q_\ell$ 层}\Term{category}{范畴}，记作
\[
\mathbb Q_\ell\text{-}\mathrm{fc}(X),
\]
是 $\mathbb Z_\ell\text{-}\mathrm{fc}(X)$ 对\Term{torsion}{挠} $\mathbb Z_\ell$ \Term{sheaf}{层}构成的\Term{thick-subcategory}{厚}\Term{abelian-category}{阿贝尔子范畴}的\Term{quotient-category}{商范畴}（参见 Gabriel 的学位论文）。$\mathbb Q_\ell\text{-}\mathrm{fc}(X)$ 的一个\Term{object}{对象}称为一个\Term{constructible-ql-sheaf}{可构造 $\mathbb Q_\ell$ 层}。
\end{SGAresult}
\Term{category}{范畴} $\mathbb Q_\ell\text{-}\mathrm{fc}(X)$ 有如下简单描述。它与 $\mathbb Z_\ell\text{-}\mathrm{fc}(X)$ 有相同的\Term{object}{对象}集；若 $F,G$ 是两个这样的\Term{object}{对象}，则
\SourcePage{278}
\[
\operatorname{Hom}_{\mathbb Q_\ell}(F,G)
=\mathbb Q_\ell\otimes_{\mathbb Z_\ell}\operatorname{Hom}_{\mathbb Z_\ell}(F,G),
\]
其中 $\operatorname{Hom}_{\mathbb Z_\ell}(F,G)$ 和 $\operatorname{Hom}_{\mathbb Q_\ell}(F,G)$ 分别表示 $\mathbb Z_\ell\text{-}\mathrm{fc}(X)$ 和 $\mathbb Q_\ell\text{-}\mathrm{fc}(X)$ 中的\Term{morphism}{态射}。事实上，若
\[
u:\mathbb Z_\ell\text{-}\mathrm{fc}(X)\longrightarrow D
\]
是从 $\mathbb Z_\ell\text{-}\mathrm{fc}(X)$ 到一个\Term{abelian-category}{阿贝尔范畴} $D$ 的\Term{exact-functor}{正合函子}，则说 $u$ 使乘以 $\mathbb Z_\ell$ 元素的映射变得\Term{invertible}{可逆}，\Term{equivalence}{等价}于说它在\Term{torsion}{挠对象}上为零。由此立即可见，以上用\Term{object}{对象}和\Term{morphism}{态射}描述的（\Term{abelian-category}{阿贝尔}）\Term{category}{范畴}解决了取商的\Term{universal-problem}{泛问题}。
若一个\Term{constructible-ql-sheaf}{可构造 $\mathbb Q_\ell$ 层}是某个\Term{constructible-zl-sheaf}{可构造}\Term{twisted-constant}{扭常值 $\mathbb Z_\ell$ 层}的像，则称它是\Term{twisted-constant}{扭常值}的。由\Term{constructible-ql-sheaf}{可构造}\Term{twisted-constant}{扭常值 $\mathbb Q_\ell$ 层}在 $\mathbb Q_\ell\text{-}\mathrm{fc}(X)$ 中生成的\Term{full-subcategory}{满子范畴}，称为\Term{constructible-ql-sheaf}{可构造}\Term{twisted-constant}{扭常值 $\mathbb Q_\ell$ 层}\Term{category}{范畴}。
由（1.2.4），若 $X$ \Term{connected}{连通}，则\Term{constructible-ql-sheaf}{可构造}\Term{twisted-constant}{扭常值 $\mathbb Q_\ell$ 层}\Term{category-equivalence}{范畴等价}于带有 $\pi_1(X)$ 的\Term{continuous-action}{连续作用}的\Term{finite-dimensional}{有限维} $\mathbb Q_\ell$ \Term{vector-space}{向量空间} $V$ 的\Term{category}{范畴}。所谓\Term{continuous-action}{连续作用}，是指一个\Term{continuous-map}{连续映射}
\[
u:\pi_1(X)\times V\longrightarrow V,
\]
满足（1.2.4）所列性质，特别地，这蕴含对每个 $g\in\pi_1(X)$，$u(g,v)$ 关于 $v$ 是 $\mathbb Q_\ell$ \Term{linear}{线性}的。
立即可见，这\Term{equivalence}{等价}于给定一个\Term{topological-group}{拓扑群}\Term{morphism}{态射} $\pi_1(X)\to\mathrm{GL}(V)$，后一个群配以其自然的 $\ell$ 进解析群\Term{topology}{拓扑}。
\SourcePage{279}
\textbf{1.4.3。\Term{constructible}{可构造} $L$ 层（$L$ 为\Term{finite}{有限} $\mathbb Q_\ell$ \Term{algebra}{代数}）。}
设 $L$ 是一个\Term{finite}{有限} $\mathbb Q_\ell$ \Term{algebra}{代数}。仿照（1.4.1）中从\Term{constructible-zl-sheaf}{可构造 $\mathbb Z_\ell$ 层}出发所作的构造，由\Term{constructible-ql-sheaf}{可构造 $\mathbb Q_\ell$ 层}\Term{category}{范畴}定义\Term{constructible-l-sheaf}{可构造 $L$ 层}\Term{category}{范畴}，继而定义\Term{constructible-l-sheaf}{可构造}\Term{twisted-constant}{扭常值 $L$ 层}\Term{category}{范畴}。若 $X$ \Term{connected}{连通}，则\Term{constructible-l-sheaf}{可构造}\Term{twisted-constant}{扭常值 $L$ 层}\Term{category-equivalence}{范畴等价}于带有 $\pi_1(X)$ 的\Term{continuous-action}{连续作用}的\Term{finite-dimensional}{有限维} $L$ \Term{vector-space}{向量空间}\Term{category}{范畴}；这里的\Term{action}{作用}是指 $\pi_1(X)$ 在底层 $\mathbb Q_\ell$ \Term{vector-space}{向量空间}上的一个\Term{continuous-action}{连续作用}，且它与 $L$ 的\Term{action}{作用}交换。

\textbf{1.4.4。} 若 $A$ 是一个交换\Term{finite}{有限} $\mathbb Z_\ell$ \Term{algebra}{代数}，$F,G$ 是两个\Term{constructible-a-sheaf}{可构造 $A$ 层}，则\Term{projective-system}{投射系}
\[
F\otimes_A G=(F_n\otimes_{A_X}G_n)_{n\in\mathbb N},
\]
配以显然的\Term{transition-morphism}{过渡态射}，是一个 $\mathbb Z_\ell$ \Term{sheaf}{层}。此外，$A$ 显然\Term{action}{作用}于它的每个\Term{set-component}{分量}，因而可以自然地赋予它一个 $A$ \Term{sheaf}{层}结构。这样定义了一个称为\Term{tensor-product}{张量积}的\Term{right-exact-functor}{右正合}\Term{bifunctor}{双函子}：
\[
A\text{-}\mathrm{fc}(X)\times A\text{-}\mathrm{fc}(X)\longrightarrow A\text{-}\mathrm{fc}(X).
\]
现设 $F,G$ 是两个 $A$ \Term{sheaf}{层}，且假设 $F$ \Term{locally-free}{局部自由}，即对每个整数 $n$，$F_n$ 是一个\Term{locally-free}{局部自由}的 $(A/\ell^{n+1}A)_X$ 模。于是，与（1.3.2）中一样，可知\Term{projective-system}{投射系}
\SourcePage{280}
\[
\mathcal{H}om_A(F,G)=\bigl(\mathcal{H}om_{A_n}(F_n,G_n)\bigr)_{n\in\mathbb N},
\]
配以显然的\Term{transition-morphism}{过渡态射}，是一个 $\mathbb Z_\ell$ \Term{sheaf}{层}。此外，$A$ 自然地\Term{action}{作用}于它的各\Term{set-component}{分量}，因而它\Term{canonical}{典范}地带有 $A$ \Term{sheaf}{层}结构。
与（1.3.3）中一样，可以定义\Term{invertible-a-sheaf}{可逆 $A$ 层}的概念；若 $L$ 是这样的一个 $A$ \Term{sheaf}{层}，则对每个自然数 $r$，赋予表达式 $L^{\otimes r}$ 以意义。（1.3.3）的内容可毫无困难地推广到这一情形。

\textbf{1.4.5。} 设 $F$ 是一个\Term{constructible-zl-sheaf}{可构造 $\mathbb Z_\ell$ 层}，以 $p$ 表示\Term{canonical-functor}{典范函子} $\mathbb Z_\ell\text{-}\mathrm{fc}(X)\to\mathbb Q_\ell\text{-}\mathrm{fc}(X)$。把\Term{functor}{函子} $G\mapsto F\otimes G$ 与 $p$ 复合，所得\Term{functor}{函子}使乘以 $\mathbb Z_\ell$ 元素的映射变得\Term{invertible}{可逆}，因此定义了一个\Term{functor}{函子}
\[
\mathbb Q_\ell\text{-}\mathrm{fc}(X)\longrightarrow\mathbb Q_\ell\text{-}\mathrm{fc}(X),
\]
它是\Term{right-exact-functor}{右正合}的。对第一个变量也应用同样的论证，就定义了一个称为\Term{tensor-product}{张量积}的\Term{right-exact-functor}{右正合}\Term{bifunctor}{双函子}：
\[
\mathbb Q_\ell\text{-}\mathrm{fc}(X)\times\mathbb Q_\ell\text{-}\mathrm{fc}(X)
\longrightarrow\mathbb Q_\ell\text{-}\mathrm{fc}(X).
\]
若 $X$ \Term{connected}{连通}，$F,G$ 是两个\Term{constructible-ql-sheaf}{可构造}\Term{twisted-constant}{扭常值 $\mathbb Q_\ell$ 层}，分别对应于 $\pi_1(X)$ 在\Term{vector-space}{向量空间} $V,W$ 上的表示，则\Term{tensor-product}{张量积} $F\otimes G$ 仍是\Term{locally-constant}{局部常值}的，并对应于 $\pi_1(X)$ 在 $V\otimes W$ 上的对角表示。
\SourcePage{281}
\textbf{1.4.6。} 设 $F$ 是一个\Term{locally-free}{局部自由}的\Term{constructible-zl-sheaf}{可构造 $\mathbb Z_\ell$ 层}（1.3.3）。把\Term{functor}{函子} $G\mapsto\mathcal{H}om(F,G)$（1.3.2）与 $p$ 复合，所得\Term{functor}{函子}使乘以 $\mathbb Z_\ell$ 元素的映射变得\Term{invertible}{可逆}，因此定义了一个\Term{functor}{函子}
\[
\mathcal{H}om(F,\cdot):\mathbb Q_\ell\text{-}\mathrm{fc}(X)
\longrightarrow\mathbb Q_\ell\text{-}\mathrm{fc}(X).
\]
若 $F,F'$ 是两个\Term{locally-free}{局部自由}的\Term{constructible-zl-sheaf}{可构造 $\mathbb Z_\ell$ 层}，且它们定义\Term{isomorphism}{同构}的 $\mathbb Q_\ell$ \Term{sheaf}{层}，则存在一个 $\mathbb Z_\ell$ \Term{sheaf}{层}的\Term{morphism}{态射}
\[
u:F\longrightarrow F',
\]
其\Term{kernel}{核}与\Term{cokernel}{余核}都是\Term{torsion}{挠对象}。由此立刻可知，从 $\mathbb Q_\ell\text{-}\mathrm{fc}(X)$ 到自身的\Term{functor}{函子} $\mathcal{H}om(F,\cdot)$ 和 $\mathcal{H}om(F',\cdot)$ \Term{isomorphism}{同构}。
现设 $F$ 是一个\Term{constructible-ql-sheaf}{可构造}\Term{twisted-constant}{扭常值 $\mathbb Q_\ell$ 层}。仍以 $F$ 表示相应的 $\mathbb Z_\ell$ \Term{sheaf}{层}。把 $F$ 在 $X$ 的各\Term{connected-component}{连通分支} $Y$ 上的\Term{restriction}{限制}解释为\Term{continuous}{连续}的 $\pi_1(Y)$ 模，可知 $F$ 在\Term{isomorphism}{同构}意义下唯一地写为
\[
F=P\oplus T,
\]
其中 $P$ \Term{locally-free}{局部自由}，$T$ 为\Term{torsion}{挠对象}。这使我们能够对每个 $\mathbb Q_\ell$ \Term{sheaf}{层} $G$，定义一个 $\mathbb Q_\ell$ \Term{sheaf}{层}
\[
\mathcal{H}om(F,G)=\mathcal{H}om(P,G).
\]
$\mathbb Q_\ell$ \Term{sheaf}{层}的一个\Term{isomorphism}{同构} $u:F\xrightarrow{\sim}F'$，对每个 $G$ 定义一个\Term{isomorphism}{同构}
\[
\mathcal{H}om(F',G)\xrightarrow{\sim}\mathcal{H}om(F,G).
\]
更确切地说，这样定义了一个\Term{bifunctor}{双函子}
\[
\mathcal{H}om:(\mathbb Q_\ell\text{-}\mathrm{fct}(X))^\circ\times
\mathbb Q_\ell\text{-}\mathrm{fc}(X)\longrightarrow\mathbb Q_\ell\text{-}\mathrm{fc}(X),
\]
\SourcePage{282}
其中 $\mathbb Q_\ell\text{-}\mathrm{fct}(X)$ 表示\Term{constructible-ql-sheaf}{可构造}\Term{twisted-constant}{扭常值 $\mathbb Q_\ell$ 层}的\Term{category}{范畴}。此\Term{bifunctor}{双函子}具有通常的\Term{exactness}{正合性质}。

\textbf{1.4.7。} 现设 $L$ 是 $\mathbb Q_\ell$ 的一个\Term{finite-type}{有限型}\Term{group-extension}{扩张}，以 $A$ 表示 $\mathbb Z_\ell$ 在 $L$ 中的\Term{integral-closure}{整闭包}；于是 $A$ 是一个\Term{finite-type}{有限型}自由 $\mathbb Z_\ell$ \Term{algebra}{代数}，也是一个 \Term{dedekind-ring}{Dedekind 环}。\Term{constructible-l-sheaf}{可构造 $L$ 层}\Term{category}{范畴}\Term{isomorphism}{同构}于\Term{constructible-a-sheaf}{可构造 $A$ 层}\Term{category}{范畴}对\Term{constructible}{可构造}\Term{torsion}{挠} $A$ \Term{sheaf}{层}\Term{category}{范畴}的\Term{quotient-category}{商范畴}。这使得我们能够与（1.4.6）一样，由（1.4.4）的\Term{tensor-product}{张量积}定义一个\Term{tensor-product}{张量积}
\[
L\text{-}\mathrm{fc}(X)\times L\text{-}\mathrm{fc}(X)\longrightarrow L\text{-}\mathrm{fc}(X).
\]
由于环 $A$ 是 \Term{dedekind-ring}{Dedekind 环}，（1.4.6）中的论证表明，每个\Term{twisted-constant}{扭常值 $A$ 层}都是一个\Term{locally-free}{局部自由} $A$ \Term{sheaf}{层}与一个\Term{torsion}{挠} $A$ \Term{sheaf}{层}的\Term{direct-sum}{直和}。于是，以 $L\text{-}\mathrm{fct}(X)$ 表示\Term{constructible-l-sheaf}{可构造}\Term{twisted-constant}{扭常值 $L$ 层}\Term{category}{范畴}，便可与（1.4.6）一样定义一个\Term{bifunctor}{双函子}
\[
\mathcal{H}om:(L\text{-}\mathrm{fct}(X))^\circ\times L\text{-}\mathrm{fc}(X)
\longrightarrow L\text{-}\mathrm{fc}(X),
\]
它满足通常的\Term{exactness}{正合性质}。
\SGAsection{1.5}{可构造 $(\mathrm{AR},\mathbb Z_\ell)$ 层}
\begin{SGAresult}{定义}{1.5.1}
设 $X$ 是一个\Term{locally-noetherian}{局部诺特}\Term{prescheme}{预概形}。$\operatorname{Hom}(\mathbb N^\circ,\mathrm{Ab}(X))$ 的一个\Term{object}{对象} $F$ 称为\Term{constructible-ar-zl-sheaf}{可构造 $(\mathrm{AR},\mathbb Z_\ell)$ 层}，若 $X$ 有一个由\Term{noetherian}{诺特}\Term{open-subset}{开集}组成的\Term{covering}{覆盖}
\[
(U_i)_{i\in I},
\]
\SourcePage{283}
使 $F$ 在每个 $U_i$ 上的\Term{restriction}{限制}都是一个 $(\mathrm{AR},\ell\mathbb Z)$ \Term{noetherian}{诺特}\Term{object}{对象}（5.1.3）。
\end{SGAresult}
由（V 5.2.1）立即可知，由\Term{constructible-ar-zl-sheaf}{可构造 $(\mathrm{AR},\mathbb Z_\ell)$ 层}在 $\operatorname{Hom}(\mathbb N^\circ,\mathrm{Ab}(X))$ 中生成的\Term{full-subcategory}{满子范畴}对于取\Term{kernel}{核}和\Term{cokernel}{余核}\Term{stable}{稳定}。特别地，它是一个\Term{abelian-category}{阿贝尔范畴}。

\textbf{1.5.2。\Term{locally-ar-null}{局部 AR 零}的\Term{constructible}{可构造} $(\mathrm{AR},\mathbb Z_\ell)$ 层。} 显然，$\operatorname{Hom}(\mathbb N^\circ,\mathrm{Ab}(X))$ 的一个 \Term{ar-null}{AR 零}\Term{object}{对象}（V 2.2.1）是\Term{constructible-ar-zl-sheaf}{可构造 $(\mathrm{AR},\mathbb Z_\ell)$ 层}，当且仅当其各\Term{set-component}{分量}\Term{constructible}{可构造}。由此，特别根据（V 2.2.2）与（SGA 4 IX 2.6），\Term{ar-null}{AR 零}的\Term{constructible-ar-zl-sheaf}{可构造 $(\mathrm{AR},\mathbb Z_\ell)$ 层}\Term{category}{范畴}，是由\Term{constructible-ar-zl-sheaf}{可构造 $(\mathrm{AR},\mathbb Z_\ell)$ 层}在 $\operatorname{Hom}(\mathbb N^\circ,\mathrm{Ab}(X))$ 中生成的\Term{full-subcategory}{满子范畴}的一个\Term{thick-subcategory}{厚}\Term{abelian-category}{阿贝尔子范畴}。
约定称 $\operatorname{Hom}(\mathbb N^\circ,\mathrm{Ab}(X))$ 的一个\Term{object}{对象} $P$ 是\Term{locally-ar-null}{局部 AR 零}的，若 $X$ 有一个对于\Term{etale-topology}{平展拓扑}（或\Term{equivalence}{等价}地，对于 Zariski \Term{topology}{拓扑}）的开\Term{covering}{覆盖}，使 $P$ 在其中每个\Term{open-subset}{开集}上都是 \Term{ar-null}{AR 零}的。通过\Term{localization}{局部化}归约到前一情形，可知\Term{locally-ar-null}{局部 AR 零}的\Term{constructible-ar-zl-sheaf}{可构造 $(\mathrm{AR},\mathbb Z_\ell)$ 层}\Term{category}{范畴}，是\Term{constructible-ar-zl-sheaf}{可构造 $(\mathrm{AR},\mathbb Z_\ell)$ 层}\Term{category}{范畴}的一个\Term{thick-subcategory}{厚}\Term{abelian-category}{阿贝尔子范畴}。
注意，若 $X$ \Term{noetherian}{诺特}，则\Term{constructible-ar-zl-sheaf}{可构造 $(\mathrm{AR},\mathbb Z_\ell)$ 层}恰是 $\operatorname{Hom}(\mathbb N^\circ,\mathrm{Ab}(X))$ 中的\Term{noetherian}{诺特} $(\mathrm{AR},\ell\mathbb Z)$ 进\Term{object}{对象}，而\Term{locally-ar-null}{局部 AR 零}\Term{equivalence}{等价}于 \Term{ar-null}{AR 零}。
\begin{SGAresult}{定义}{1.5.3}
设 $X$ 是一个\Term{locally-noetherian}{局部诺特}\Term{prescheme}{预概形}。$X$ 上的\Term{constructible-ar-zl-sheaf}{可构造 $(\mathrm{AR},\mathbb Z_\ell)$ 层}\Term{category}{范畴}，记作
\SourcePage{284}
\[
(\mathrm{AR},\mathbb Z_\ell)\text{-}\mathrm{fc}(X),
\]
是以下\Term{abelian-category}{阿贝尔}\Term{quotient-category}{商范畴}：把由\Term{constructible-ar-zl-sheaf}{可构造 $(\mathrm{AR},\mathbb Z_\ell)$ 层}在
\[
\operatorname{Hom}(\mathbb N^\circ,\mathrm{Ab}(X))
\]
中生成的\Term{full-subcategory}{满子范畴}，对其中由\Term{locally-ar-null}{局部 AR 零}的\Term{constructible-ar-zl-sheaf}{可构造 $(\mathrm{AR},\mathbb Z_\ell)$ 层}生成的\Term{full-subcategory}{满}\Term{thick-subcategory}{厚子范畴}取商。
\end{SGAresult}
\textbf{记号 1.5.4。} 有需要时，把由\Term{constructible-ar-zl-sheaf}{可构造 $(\mathrm{AR},\mathbb Z_\ell)$ 层}在 $\operatorname{Hom}(\mathbb N^\circ,\mathrm{Ab}(X))$ 中生成的\Term{full-subcategory}{满子范畴}记作
\[
\mathrm U(X).
\]
\Term{constructible-zl-sheaf}{可构造 $\mathbb Z_\ell$ 层}特别也是\Term{constructible-ar-zl-sheaf}{可构造 $(\mathrm{AR},\mathbb Z_\ell)$ 层}，因此显然定义了一个\Term{functor}{函子}
\[
u:\mathbb Z_\ell\text{-}\mathrm{fc}(X)\longrightarrow
(\mathrm{AR},\mathbb Z_\ell)\text{-}\mathrm{fc}(X).
\]
\begin{SGAresult}{命题}{1.5.5}
上述\Term{functor}{函子} $u$ 是一个\Term{equivalence-of-categories}{范畴等价}。
\end{SGAresult}
\begin{proof}
为证 $u$ \Term{fully-faithful}{全忠实}且\Term{essentially-surjective}{本质满射}，可“通过\Term{gluing}{粘合}”（如命题 1.1.3 的证明中所作的）归约到相应的断言（V 5.1.3）。
\end{proof}
\SourcePage{285}
\SGAsection{2}{$\ell$ 进上同调的形式体系}
\textbf{2.1。\Term{inverse-image}{逆像}。} 设 $f:X\to Y$ 是\Term{locally-noetherian}{局部诺特}\Term{prescheme}{预概形}的一个\Term{morphism}{态射}。\Term{inverse-image}{逆像}\Term{functor}{函子}
\[
f^*:\mathrm{Ab}(Y)\longrightarrow\mathrm{Ab}(X)
\]
是\Term{exact-functor}{正合}的，因而显然延拓为一个同样记号的\Term{exact-functor}{正合函子}
\[
f^*:\operatorname{Hom}(\mathbb N^\circ,\mathrm{Ab}(Y))\longrightarrow
\operatorname{Hom}(\mathbb N^\circ,\mathrm{Ab}(X)).
\]
后一个\Term{functor}{函子}显然把\Term{constructible-zl-sheaf}{可构造 $\mathbb Z_\ell$ 层}变为\Term{constructible-zl-sheaf}{可构造 $\mathbb Z_\ell$ 层}，把\Term{locally-ar-null}{局部 AR 零}\Term{projective-system}{投射系}变为\Term{locally-ar-null}{局部 AR 零}\Term{projective-system}{投射系}。采用（1.5.4）的记号，它因此诱导两个“\Term{inverse-image}{逆像}”\Term{functor}{函子}
\[
\begin{aligned}
f^*&:\mathbb Z_\ell\text{-}\mathrm{fc}(Y)\longrightarrow\mathbb Z_\ell\text{-}\mathrm{fc}(X),\\
f^*&:\mathrm U(Y)\longrightarrow\mathrm U(X),
\end{aligned}
\]
第二个是\Term{exact-functor}{正合}的。
由\Term{equivalence-of-categories}{范畴等价}（1.5.5），以及 $f^*$ 把\Term{locally-ar-null}{局部 AR 零}\Term{projective-system}{投射系}变为\Term{locally-ar-null}{局部 AR 零}\Term{projective-system}{投射系}这一事实，便知\Term{inverse-image}{逆像}\Term{functor}{函子}
\[
f^*:\mathbb Z_\ell\text{-}\mathrm{fc}(Y)\longrightarrow\mathbb Z_\ell\text{-}\mathrm{fc}(X)
\]
也是\Term{exact-functor}{正合}的。
若 $g:Y\to Z$ 是\Term{locally-noetherian}{局部诺特}\Term{prescheme}{预概形}的另一个\Term{morphism}{态射}，则\Term{functor}{函子} $f^*$ 与 $g^*$ 显然通过\Term{isomorphism}{同构}
\[
(g\circ f)^*\xrightarrow{\sim}f^*\circ g^*
\]
联系起来，并满足通常的“结合”性质（SGA 1 VI 7 及以下各节）。
\SourcePage{286}
\textbf{2.2。\Term{higher-direct-image}{高阶直像}。} 设
\[
f:X\longrightarrow Y
\]
是\Term{locally-noetherian}{局部诺特}\Term{prescheme}{预概形}之间的一个\Term{separated-morphism}{分离}\Term{finite-type}{有限型}\Term{morphism}{态射}。回顾\Term{etale-cohomology}{平展上同调}的下述结果（SGA 4 XVII；\Term{proper}{固有}情形见 SGA 4 XIV 1.1）。
\begin{SGAresult}{\Term{lemma}{引理}}{2.2.1}
设 $A$ 是一个被非零整数 $n$ \Term{annihilate}{零化}的\Term{commutative-ring}{交换环}，$M$ 是一个\Term{constructible}{可构造} $(A,\mathrm{Ab}(X))$ 模。对每个整数 $i$，\Term{sheaf}{层}
\[
\mathrm R^i_{!}f(M)
\]
是一个\Term{constructible}{可构造} $(A,\mathrm{Ab}(Y))$ 模。
\end{SGAresult}
把上面的结论\Term{restriction}{限制}到 $Y$ 的\Term{noetherian}{诺特}\Term{open-subset}{开集}，再应用（V 5.3.1），便得到下述\Term{lemma}{引理}。
\begin{SGAresult}{\Term{lemma}{引理}}{2.2.2}
对每个整数 $i$ 和 $X$ 上的每个\Term{constructible-ar-zl-sheaf}{可构造 $(\mathrm{AR},\mathbb Z_\ell)$ 层} $F$，\Term{projective-system}{投射系}
\[
\mathrm R^i_{!}f(F)
\]
是一个\Term{constructible-ar-zl-sheaf}{可构造 $(\mathrm{AR},\mathbb Z_\ell)$ 层}。
\end{SGAresult}
由此，\Term{exact-delta-functor}{正合 $\partial$ 函子}
\[
(\mathrm R^i_{!}f)_{i\in\mathbb Z}:\mathrm{Ab}(X)\longrightarrow\mathrm{Ab}(Y)
\]
显然延拓为一个同样记号的\Term{exact-delta-functor}{正合 $\partial$ 函子}
\[
\operatorname{Hom}(\mathbb N^\circ,\mathrm{Ab}(X))\longrightarrow
\operatorname{Hom}(\mathbb N^\circ,\mathrm{Ab}(Y)),
\]
其\Term{restriction}{限制}给出一个\Term{exact-delta-functor}{正合 $\partial$ 函子}
\[
\mathrm U(X)\longrightarrow\mathrm U(Y).\tag{1}
\]
另一方面，各\Term{functor}{函子} $\mathrm R^i_{!}f$ 把\Term{locally-ar-null}{局部 AR 零}\Term{projective-system}{投射系}变为\Term{locally-ar-null}{局部 AR 零}\Term{projective-system}{投射系}，因而取商以后，$\partial$ \Term{functor}{函子} (1) 给出一个\Term{exact-delta-functor}{正合 $\partial$ 函子}
\SourcePage{287}
\[
(\mathrm R^i_{!}f)_{i\in\mathbb Z}:
(\mathrm{AR},\mathbb Z_\ell)\text{-}\mathrm{fc}(X)\longrightarrow
(\mathrm{AR},\mathbb Z_\ell)\text{-}\mathrm{fc}(Y).
\]
再利用\Term{equivalence-of-categories}{范畴等价}（1.5.5），后一个 $\partial$ \Term{functor}{函子}便解释为\Term{exact-delta-functor}{正合 $\partial$ 函子}
\[
\mathbb Z_\ell\text{-}\mathrm{fc}(X)\longrightarrow\mathbb Z_\ell\text{-}\mathrm{fc}(Y).
\]
\begin{SGAresult}{定义}{2.2.3}
\Term{functor}{函子}
\[
\mathrm R^i_{!}f:\mathbb Z_\ell\text{-}\mathrm{fc}(X)\longrightarrow\mathbb Z_\ell\text{-}\mathrm{fc}(Y)
\]
称为 $f$ 的第 $i$ 个\Term{direct-image-proper-support}{固有支集直像}\Term{functor}{函子}。
\end{SGAresult}
\textbf{注记。} 假设在 $X,Y$ 为\Term{excellent}{优良}时，已经对普通\Term{higher-direct-image}{高阶直像}证明了\Term{lemma}{引理}（2.2.1）的相应结论（\Term{characteristic-zero}{特征零}的情形已在 SGA 4 XIX 5.1 中证明），那么在这一情形也可类似地定义 $\ell$ 进\Term{cohomology}{上同调}中的\Term{higher-direct-image}{高阶直像}。下一小节的结果也可以在这个框架中展开。
\begin{SGAresult}{命题}{2.2.4}
\textbf{（\Term{leray-spectral-sequence}{Leray 谱序列}。）} 设 $f:X\to Y$、$g:Y\to Z$ 是\Term{locally-noetherian}{局部诺特}\Term{prescheme}{预概形}之间的两个\Term{separated-morphism}{分离}\Term{finite-type}{有限型}\Term{morphism}{态射}，$F$ 是一个\Term{constructible-zl-sheaf}{可构造 $\mathbb Z_\ell$ 层}（相应地，一个\Term{constructible-ar-zl-sheaf}{可构造 $(\mathrm{AR},\mathbb Z_\ell)$ 层}）。在 $\mathbb Z_\ell\text{-}\mathrm{fc}(Z)$ 中存在一个关于 $F$ \Term{functorial}{函子性}的\Term{biregular-spectral-sequence}{双正则谱序列}（EGA $0_{\mathrm{III}}$ 11.1.3）：
\[
\mathrm R^i_{!}g\bigl(\mathrm R^j_{!}f(F)\bigr)
\Longrightarrow\mathrm R^{i+j}_{!}(g\circ f)(F).
\]
\end{SGAresult}
\begin{proof}
由\Term{equivalence-of-categories}{范畴等价}（1.5.5）及\Term{constructible-zl-sheaf}{可构造 $\mathbb Z_\ell$ 层}的\Term{higher-direct-image}{高阶直像}的定义方式，只须对\Term{constructible-ar-zl-sheaf}{可构造 $(\mathrm{AR},\mathbb Z_\ell)$ 层}证明相应断言。对任意\Term{abelian-category}{阿贝尔范畴} $C$，约定以
\[
\operatorname{Sp}(C)
\]
\SourcePage{288}
表示在\Term{degree}{次数} $2$ 开始、取值于 $C$ 的\Term{spectral-sequence}{谱序列}\Term{category}{范畴}。对 $\operatorname{Hom}(\mathbb N^\circ,\mathrm{Ab}(X))$ 的每个\Term{object}{对象} $F$，令其对应于把普通 Leray \Term{spectral-sequence}{谱序列}延拓到\Term{projective-system}{投射系}而得的\Term{spectral-sequence}{谱序列}
\[
\mathrm R^i_{!}g\bigl(\mathrm R^j_{!}f(F)\bigr)
\Longrightarrow\mathrm R^{i+j}_{!}(g\circ f)(F).\tag{S}
\]
显然，这给出一个\Term{functor}{函子}
\[
\operatorname{Hom}(\mathbb N^\circ,\mathrm{Ab}(X))\longrightarrow
\operatorname{Sp}\bigl(\operatorname{Hom}(\mathbb N^\circ,\mathrm{Ab}(Z))\bigr).
\]
根据（2.2.1），以及 $\mathrm U(Z)$ 在 $\operatorname{Hom}(\mathbb N^\circ,\mathrm{Ab}(Z))$ 中为\Term{thick-subcategory}{厚子范畴}这一事实，取\Term{restriction}{限制}得到一个\Term{functor}{函子}
\[
\mathrm U(X)\longrightarrow\operatorname{Sp}(\mathrm U(Z)).\tag{1}
\]
\Term{canonical-functor}{典范函子} $\mathrm U(Z)\to(\mathrm{AR},\mathbb Z_\ell)\text{-}\mathrm{fc}(Z)$ 是\Term{exact-functor}{正合}的，因此由 (1) 得到一个\Term{additive-functor}{加性函子}
\[
E:\mathrm U(X)\longrightarrow
\operatorname{Sp}\bigl((\mathrm{AR},\mathbb Z_\ell)\text{-}\mathrm{fc}(Z)\bigr).\tag{2}
\]
对 $\operatorname{Hom}(\mathbb N^\circ,\mathrm{Ab}(X))$ 的每个\Term{object}{对象} $F$，所对应的\Term{spectral-sequence}{谱序列} (S) 都是\Term{biregular-spectral-sequence}{双正则}的：这是由\Term{projective-system}{投射系} $F$ 各\Term{set-component}{分量}的相应断言，以及所考虑的\Term{spectral-sequence}{谱序列}支承于第一象限这一事实得到的。因此，$\mathrm U(X)$ 的每个\Term{object}{对象}所对应的\Term{spectral-sequence}{谱序列}都是\Term{biregular-spectral-sequence}{双正则}的。
此外，若 $u:F\to G$ 是 $\mathrm U(X)$ 的一个\Term{morphism}{态射}，且其\Term{kernel}{核}与\Term{cokernel}{余核}都是\Term{locally-ar-null}{局部 AR 零}的，则由于\Term{higher-direct-image}{高阶直像}把\Term{locally-ar-null}{局部 AR 零}\Term{projective-system}{投射系}变为\Term{locally-ar-null}{局部 AR 零}\Term{projective-system}{投射系}，可知 $E(u)$ 在各项 $E_2^{pq}$ 上都是\Term{isomorphism}{同构}；而所考虑的\Term{spectral-sequence}{谱序列}都是\Term{biregular-spectral-sequence}{双正则}的，故 $E(u)$ 是一个\Term{isomorphism}{同构}。
由此，取商后 (2) 定义了一个\Term{functor}{函子}
\[
(\mathrm{AR},\mathbb Z_\ell)\text{-}\mathrm{fc}(X)\longrightarrow
\operatorname{Sp}\bigl((\mathrm{AR},\mathbb Z_\ell)\text{-}\mathrm{fc}(Z)\bigr),\tag{3}
\]
\SourcePage{289}
它把每个\Term{constructible-ar-zl-sheaf}{可构造 $(\mathrm{AR},\mathbb Z_\ell)$ 层} $F$ 明确地对应到一个\Term{biregular-spectral-sequence}{双正则谱序列}：
\[
\mathrm R^i_{!}f\bigl(\mathrm R^j_{!}g(F)\bigr)
\Longrightarrow\mathrm R^{i+j}_{!}(g\circ f)(F).
\]
\end{proof}
\textbf{2.2.3。其他性质。}

\textbf{A）} 设 $f:X\to Y$ 是\Term{locally-noetherian}{局部诺特}\Term{prescheme}{预概形}之间的一个\Term{separated-morphism}{分离}\Term{finite-type}{有限型}\Term{morphism}{态射}。假设 $f$ 的\Term{relative-dimension}{相对维数}小于或等于整数 $d$。那么，对每个\Term{constructible-ar-zl-sheaf}{可构造 $(\mathrm{AR},\mathbb Z_\ell)$ 层} $F$，都有
\[
\mathrm R^i_{!}f(F)=0\qquad\text{当 }i>2d.
\]
这由 $F$ 各\Term{set-component}{分量}的相应断言得到（SGA 4 X 4.1 和 5.2）。

\textbf{B）\Term{proper}{固有}\Term{base-change}{基变换}。}
\[
\begin{tikzcd}
X'\arrow[r,"g'"]\arrow[d,"f'"']&X\arrow[d,"f"]\\
Y'\arrow[r,"g"]&Y
\end{tikzcd}
\]
设 $f:X\to Y$、$g:Y'\to Y$ 是\Term{locally-noetherian}{局部诺特}\Term{prescheme}{预概形}的\Term{morphism}{态射}，且 $f$ 是\Term{separated-morphism}{分离}\Term{finite-type}{有限型}的。令 $X'=X\times_Y Y'$（由于 $f$ 是\Term{finite-type}{有限型}的，它是\Term{locally-noetherian}{局部诺特}的），并以
\[
f':X'\longrightarrow Y',\qquad g':X'\longrightarrow X
\]
表示\Term{canonical-morphism}{典范态射}。对每个整数 $i$ 和每个\Term{constructible-zl-sheaf}{可构造 $\mathbb Z_\ell$ 层} $F$，在 $\mathbb Z_\ell\text{-}\mathrm{fc}(Y')$ 中存在关于 $F$ \Term{functorial}{函子性}的\Term{isomorphism}{同构}
\[
w_i:g^*\bigl(\mathrm R^i_{!}f(F)\bigr)\longrightarrow
\mathrm R^i_{!}f'(g'^*F).
\]
\SourcePage{290}
此外，若 $0\to F'\to F\to F''\to0$ 是 $\mathbb Z_\ell\text{-}\mathrm{fc}(X)$（相应地，$(\mathrm{AR},\mathbb Z_\ell)\text{-}\mathrm{fc}(X)$）中的一个\Term{exact-sequence}{正合列}，则 $(w_i)_{i\in\mathbb N}$ 与边缘\Term{operator}{算子}交换，换言之，下列图表交换：
\[
\begin{tikzcd}[column sep=large]
g^*\bigl(\mathrm R^i_{!}f(F'')\bigr)\arrow[r,"w_i"]\arrow[d,"g^*(\delta)"']&
\mathrm R^i_{!}f'(g'^*F'')\arrow[d,"\delta"]\\
g^*\bigl(\mathrm R^{i+1}_{!}f(F')\bigr)\arrow[r,"w_{i+1}"]&
\mathrm R^{i+1}_{!}f'(g'^*F').
\end{tikzcd}
\]
\begin{proof}
只须对\Term{constructible-ar-zl-sheaf}{可构造 $(\mathrm{AR},\mathbb Z_\ell)$ 层}证明相应断言。而这一断言可由\Term{torsion}{挠}\Term{sheaf}{层}\Term{projective-system}{投射系}的相应断言取商得到；后者则是\Term{proper-base-change}{固有基变换定理}的直接推论（\Term{proper}{固有}情形见 SGA 4 XII 5.1，一般情形见 SGA 4 XVII）。
\end{proof}
\textbf{C）关于\Term{open-subset}{开集}及其补集的\Term{exact-sequence}{正合列}。}
设 $f:X\to Y$ 是\Term{locally-noetherian}{局部诺特}\Term{prescheme}{预概形}之间的一个\Term{compactifiable}{可紧化}\Term{morphism}{态射}，$U$ 是 $X$ 的一个\Term{open-subset}{开集}，$Z$ 是其闭补集。对每个\Term{constructible-zl-sheaf}{可构造 $\mathbb Z_\ell$ 层} $F$，在 $\mathbb Z_\ell\text{-}\mathrm{fc}(Y)$ 中存在一个关于 $F$ \Term{functorial}{函子性}的\Term{exact-sequence}{正合列}：
\[
\cdots\longrightarrow\mathrm R^i_{!}f_U(F/U)\longrightarrow\mathrm R^i_{!}f(F)
\longrightarrow\mathrm R^i_{!}f_Z(F/Z)\longrightarrow
\mathrm R^{i+1}_{!}f_U(F/U)\longrightarrow\cdots.
\]
\begin{proof}与 B）的证明类似。\end{proof}
\textbf{2.2.5。} 本节为 $\mathbb Z_\ell$ \Term{sheaf}{层}建立的\Term{form}{形式}体系，可以毫无困难地推广到 $A$ \Term{sheaf}{层}（$A$ 是一个\Term{finite}{有限} $\mathbb Z_\ell$ \Term{algebra}{代数}）及 $L$ \Term{sheaf}{层}（$L$ 是 $\mathbb Q_\ell$ 的一个\Term{finite-extension}{有限扩张}）。例如在 $\mathbb Q_\ell$ \Term{sheaf}{层}的情形，证明的关键在于：
\SourcePage{291}
对 $\mathbb Z_\ell$ \Term{sheaf}{层}所考虑的全部\Term{functor}{函子}，都把\Term{torsion}{挠} $\mathbb Z_\ell$ \Term{sheaf}{层}变为\Term{torsion}{挠} $\mathbb Z_\ell$ \Term{sheaf}{层}。
\SGAsection{3}{与一个圈相联系的 $\ell$ 进上同调类}
\textbf{3.1。} 若 $X$ 是\Term{algebraically-closed-field}{代数闭域} $k$ 上的一个\Term{proper}{固有}\Term{scheme}{概形}，或一个\Term{smooth}{光滑}\Term{finite-type}{有限型}\Term{scheme}{概形}，则我们知道，对每个与 $k$ 的\Term{characteristic}{特征}互素的\Term{constructible-sheaf}{可构造}\Term{torsion}{挠}\Term{abelian-sheaf}{阿贝尔层} $M$，各\Term{cohomology-group}{上同调群} $\mathrm H^i(X,M)$ 都是\Term{finite-length}{有限长度}的\Term{abelian-group}{阿贝尔群}。因此（2.2.1，注记），对每个 $\mathbb Z_\ell$ \Term{sheaf}{层} $F$（$\ell$ 与 $k$ 的\Term{characteristic}{特征}互素），可定义一个\Term{noetherian}{诺特}的\Term{abelian-group}{阿贝尔群} $(\mathrm{AR},\ell)$ 进\Term{projective-system}{投射系}
\[
\widehat{\mathrm H}^{\,i}(X,F),
\]
它具有通常的性质。特别地，第 2 节所用的方法可以定义\Term{cup-product}{杯积}、$k$ 上\Term{scheme}{概形}\Term{morphism}{态射}的\Term{inverse-image}{逆像}\Term{homomorphism}{同态}，以及 $k$ 上\Term{smooth}{光滑}\Term{scheme}{概形}之间\Term{proper}{固有}\Term{morphism}{态射}的\Term{direct-image}{直像}\Term{homomorphism}{同态}（\Term{torsion}{挠}\Term{sheaf}{层}情形的定义参见第 IV 讲）。
另一方面，我们知道（1.2.4），\Term{algebraically-closed-field}{代数闭域}上的 $\mathbb Z_\ell$ \Term{sheaf}{层}\Term{category}{范畴}，\Term{equivalence}{等价}于\Term{finite-type}{有限型} $\mathbb Z_\ell$ \Term{module-category}{模范畴}；这一\Term{equivalence-of-categories}{等价}由\Term{projective-limit}{投射极限}\Term{functor}{函子}给出。在实际应用中，把\Term{finite-type}{有限型} $\mathbb Z_\ell$ 模 $\mathrm H^i(X,F)$ 与相应的 $\mathbb Z_\ell$ \Term{sheaf}{层}等同，将很方便。
\SourcePage{292}
\textbf{3.2。} 设 $\ell$ 是一个不等于 $\operatorname{car}k$ 的素数，$X$ 是 $k$ 上的一个\Term{smooth}{光滑}\Term{finite-type}{有限型}\Term{scheme}{概形}。若 $Z$ 是 $X$ 上一个\Term{codimension}{余维}为 $d$ 的\Term{algebraic-cycle}{代数圈}，则第 IV 讲\footnote{参见（SGA $4\frac12$ Cycle）。}中对每个整数 $n$，已为它定义了一个\Term{cohomology-class}{上同调类}
\[
\gamma_X^{(n)}(Z)\in\mathrm H^{2d}\bigl(X,(\mu_X^{(n)})^{\otimes d}\bigr),
\]
其中 $\mu_X^{(n)}$ 表示由 $\ell^{n+1}$ 次\Term{root-of-unity}{单位根}组成的\Term{locally-free}{局部自由} $(\mathbb Z/\ell^{n+1}\mathbb Z)_X$ 模。此外，我们已经看到，若 $n'\geq n$，则 $\gamma_X^{(n)}(Y)$ 是 $\gamma_X^{(n')}(Y)$ 在 $\mathrm H^{2d}\bigl(X,(\mu_X^{(n)})^{\otimes d}\bigr)$ 中的像，对应的\Term{morphism}{态射}为
\[
\mathrm H^{2d}\bigl(X,(\mu_X^{(n')})^{\otimes d}\bigr)
\longrightarrow\mathrm H^{2d}\bigl(X,(\mu_X^{(n)})^{\otimes d}\bigr),
\]
它由取 $\ell^{n'-n}$ 次幂的映射 $\mu_X^{(n')}\to\mu_X^{(n)}$ 诱导。
于是，对 $Z$ 定义了一个记作 $\gamma_X(Z)$ 的元素：
\[
\gamma_X(Z)\in\mathrm H^{2d}(X,\mathbb Z_\ell(d)),
\]
其中 $\mathbb Z_\ell(d)$ 已在（1.3.4）中定义。更一般地，若 $Z$ 是 $X$ 的一个\Term{algebraic-cycle}{圈}，则按\Term{linear}{线性}性为它定义一个属于
\[
\bigoplus_d\mathrm H^{2d}(X,\mathbb Z_\ell(d))
\]
的 $\ell$ 进\Term{cohomology-class}{上同调类}；它实际上是与 $Z$ 相联系的、属于各
\[
\bigoplus_d\mathrm H^{2d}\bigl(X,(\mu_X^{(n)})^{\otimes d}\bigr)
\]
的\Term{cohomology-class}{上同调类}的\Term{projective-limit}{投射极限}。

\textbf{3.3。} 通过取\Term{limit}{极限}，显然上一段中的 $\ell$ 进\Term{cohomology-class}{上同调类}满足与第 IV 讲中研究的\Term{cohomology-class}{上同调类}相同的\Term{form}{形式}性质，即：

(i) 若 $f:X\to Y$ 是 $k$ 上\Term{smooth}{光滑}\Term{prescheme}{预概形}之间的一个\Term{proper}{固有}\Term{morphism}{态射}，$Z$ 是 $X$ 上的一个\Term{algebraic-cycle}{圈}，则有
\SourcePage{293}
\[
f_*(\gamma_X(Z))=\gamma_Y(f_*(Z)).
\]
\footnote{这一相容性在（SGA $4\frac12$ Cycle）中没有证明，但不难从该处的结果\Term{pushout}{推出}。}
(ii) 若 $f:X\to Y$ 是 $k$ 上\Term{smooth}{光滑}\Term{prescheme}{预概形}之间的一个\Term{morphism}{态射}，$Z$ 是 $Y$ 上的一个\Term{algebraic-cycle}{圈}，则有
\[
\gamma_X(f^*Z)=f^*(\gamma_Y(Z)).
\]
\footnote{在非超额情形，精确陈述见（SGA $4\frac12$ Cycle 2.3.8）。}
(iii) 若 $X$ 是 $k$ 上的一个\Term{smooth}{光滑}\Term{prescheme}{预概形}，$Z_1,Z_2$ 是 $X$ 上的两个\Term{algebraic-cycle}{圈}，则有
\[
\gamma_X(Z_1\cdot Z_2)=\gamma_X(Z_1)\cdot\gamma_X(Z_2).
\]
\footnotemark[\value{footnote}]
左边的点表示\Term{algebraic-cycle}{圈}的\Term{intersection-product}{交积}，右边的点表示\Term{cohomology}{上同调}中的\Term{cup-product}{杯积}。

(iv) 一个\Term{smooth}{光滑} $k$ \Term{prescheme}{预概形}上两个\Term{algebraic-equivalence}{代数等价}的\Term{algebraic-cycle}{圈}具有相同的 $\ell$ 进\Term{cohomology-class}{上同调类}。
